\clearpage
\phantomsection\addcontentsline{toc}{chapter}{课程信息与教学进度}
\noindent{\bfseries\fontsize{12}{17}\selectfont 课程信息与教学进度}\par
官方课程为MIT数学系18.712《表示论导论》(Introduction to Representation Theory), 2010年秋季本科课程, 教师Pavel Etingof. 课程目标是介绍群、李代数与结合代数的表示论, 重点以线性空间中的对称及具体例子建立理解. 每周两次讲课, 每次1.5小时. 官网页面未列具体上课日期、钟点或教室, 此处只保留周次, 不推造日历.

官方先修为18.701、18.702(代数I、II), 或18.700、18.703(线性代数与现代代数), 要求扎实的线性代数及群、环、域基础. 本中文本前九章有限群主线补充必要群论; 后部课程原题涉及更广的代数、李代数与箭图, 所需前置在相应题解中补明. 以下大纲是原课程范围, 不等于本书正文全部覆盖.

原大纲包括: 结合代数、生成元与关系、群代数、路径代数、李代数与包络代数、不可约与不可分解、Schur引理和 $\mathfrak{sl}_2$; 稠密性、半单代数、Jordan--Hölder与Krull--Schmidt分解、表示扩张; 有限群的Maschke、平方和、对偶与张量、特征标与矩阵元素正交、酉表示和特征标表; Frobenius--Schur指标、群行列式、代数整数与整除性、Burnside定理、诱导、互反、有限一般线性群、对称群、Schur--Weyl对偶及不变量理论; 箭图、子空间三元组、根系、反射函子和Gabriel定理.

\begin{longtable}{c p{.22\textwidth} p{.60\textwidth}}
\toprule 周次&官方旧讲义阅读范围&内容参考\\\midrule\endhead
1&§1.1--1.5&代数、表示、理想与商\\
2&§1.6--1.13&生成元关系、箭图、李代数、张量与对偶\\
3&§1.14--2.8&$\mathfrak{sl}_2$、一般表示论、分解定理\\
4&§2.9--3.5&一般练习、有限群完全可约、特征标正交\\
5&§3.6--4.1&酉表示、矩阵系数、特征标表与指标\\
6&§4.2--4.9&群行列式、整除性、Burnside、诱导\\
7&§4.10--4.19&互反、对称群、钩长与Schur--Weyl对偶\\
8&§4.20--4.24&Schur多项式、一般线性群表示\\
9&§4.25--5.2&Artin、半直积、箭图入门\\
10&§5.3--5.5&$D_4$、根系与Gabriel定理\\
11&§5.6--5.9&反射函子、Coxeter元素及证明\\
12&第6章&范畴、函子、伴随与阿贝尔范畴\\\bottomrule
\end{longtable}
阅读书目除旧讲义外, 官网列 William Fulton与Joe Harris, \textit{Representation Theory: A First Course}, Springer, 1991, GTM 129; Jean-Pierre Serre, \textit{Linear Representations of Finite Groups}, Springer-Verlag, 1977, GTM 42, 均使用开头部分. 这些是课程参考书目, 本项目不归档其受版权保护的书籍, 也不以它们替换公开旧讲义.

官方考核为作业75\%、带回家期末作业25\%. 作业每周四布置, 下一周四截止; 允许合作, 要求学生实际参与并理解所写内容. 作业页列11次作业, 共51道必做原题及1道选做原题1.58; 期末PDF有4道大题. 未列独立测验、期中卷、练习考试或官方参考答案. 本书将这些公开题目全部中文化纳入同一主PDF, 保留原编号; 解答是编者独立推导, 与旧讲义给出的提示分别标明. 官方页面: \href{https://ocw.mit.edu/courses/18-712-introduction-to-representation-theory-fall-2010/pages/syllabus/}{大纲}, \href{https://ocw.mit.edu/courses/18-712-introduction-to-representation-theory-fall-2010/pages/readings/}{阅读进度}, \href{https://ocw.mit.edu/courses/18-712-introduction-to-representation-theory-fall-2010/pages/assignments/}{作业与期末}.
